% §5 Tool-Using Agents: The ReAct Extension
We extend the formal pipeline model to cover tool-using agents that execute multiple
reasoning-action cycles before producing a response.

\subsection{The \texorpdfstring{$k$}{k}-Step ReAct Agent}

A $k$-step ReAct agent \citep{yao2023react} executes $k$ interleaved reasoning steps
and tool calls before generating its final textual response.
In the voice cascade setting, the full pipeline for a $k$-step agent is:
\begin{equation}
  \text{ASR} \to \bigl[\text{LM}_1 \to \text{Tool}_1 \to \cdots
  \to \text{LM}_{k-1} \to \text{Tool}_{k-1}\bigr]
  \to \text{LM}_k \to \text{TTS}
\end{equation}
where $\text{LM}_j$ denotes the $j$-th LM inference step and $\text{Tool}_j$ denotes
the $j$-th external tool call.
The user perceives silence for the entire duration of this chain.

\subsection{Compounding Latency Theorem}

\begin{theorem}[$k$-Step Silence Gap]
\label{thm:react}
Let $\Lambda_k$ denote the silence gap of a $k$-step ReAct agent in a cascade pipeline.
Let $\lambda_{lm}$ denote the latency of a single LM inference step
(prefill plus decode to the first token of the response or tool call),
and let $\lambda_{tool}$ denote the round-trip latency of a single tool call.
We assume all $k$ LM inference steps have equal latency $\lambda_{lm}$ (in practice, later steps incur higher prefill cost as context grows with tool observations).
For intermediate ReAct steps, only time-to-first-token (TTFT) determines when the tool call is issued, while the final LM step's full decode time feeds into TTS; this formulation treats $\lambda_{lm}$ uniformly across all steps as a simplifying assumption.
Then:
\begin{equation}
  \Lambda_k = \lambda_{asr} + k \cdot \lambda_{lm} + (k-1) \cdot \lambda_{tool} + \lambda_{tts}
\end{equation}
and the irreducible floor is:
\begin{equation}
  \Lambda_k^* = \lambda_{asr}^{irr} + k \cdot \lambda_{lm}^{irr} + (k-1) \cdot \lambda_{tool}^{irr} + \lambda_{tts}^{irr}
\end{equation}
For in-process local tool calls, $\lambda_{tool}^{irr} \approx 0$ and the simplified form $\Lambda_k^* \approx \lambda_{asr}^{irr} + k \cdot \lambda_{lm}^{irr} + \lambda_{tts}^{irr}$ applies; for remote API tools, $\lambda_{tool}^{irr}$ reflects minimum network round-trip latency.
\end{theorem}

\begin{proof}
The $k$-step ReAct pipeline forms a linear dependency chain:
each LM step cannot begin until the preceding tool call (or the ASR stage, for $\text{LM}_1$)
has completed, and each tool call cannot begin until the preceding LM step has completed.
By Definition~\ref{def:decomposition}, the latency of each stage decomposes into
reducible and irreducible components.
Summing over the critical path of the chain, the total silence gap is:
\[
  \Lambda_k = \lambda_{asr} + \sum_{j=1}^{k} \lambda_{lm} + \sum_{j=1}^{k-1} \lambda_{tool}
              + \lambda_{tts}
  = \lambda_{asr} + k \cdot \lambda_{lm} + (k-1) \cdot \lambda_{tool} + \lambda_{tts}.
\]
For the irreducible floor, we apply the argument of Theorem~\ref{thm:latency_floor}:
the reducible components of each LM step and each tool call can be driven to zero
by faster models and faster infrastructure, but the irreducible component $\lambda_{lm}^{irr}$
persists for each of the $k$ LM steps.
For tool calls, $\lambda_{tool}^{irr}$ reflects the minimum service round-trip latency;
for remote APIs this is positive, but even in the best case
(local in-process tools, $\lambda_{tool}^{irr} \approx 0$) it cannot be negative.
Applying Definition~\ref{def:decomposition} to each stage and collecting
irreducible terms along the critical path yields $\Lambda_k^*$ as stated.
\end{proof}

Figure~\ref{fig:react_loop} illustrates the compounding latency structure of the
$k$-step ReAct pipeline.

\begin{figure}[t]
  \centering
  \resizebox{\linewidth}{!}{%
    % Fig 3: ReAct loop DAG showing compounding latency for k tool calls
\begin{tikzpicture}[
  node distance=0.25cm and 0.5cm,
  box/.style={rectangle, rounded corners=3pt, draw=black, thick,
              minimum width=1.4cm, minimum height=0.6cm,
              font=\scriptsize, align=center},
  tool/.style={box, fill=orange!20},
  lm/.style={box, fill=blue!15},
  arr/.style={-{Stealth[length=4pt]}, thick},
  annot/.style={font=\scriptsize}
]

% Nodes (left to right)
\node[box, fill=gray!15] (asr)  {ASR};
\node[lm,  right=of asr]  (lm1)  {LM~1\\(reason)};
\node[tool, right=of lm1] (t1)   {Tool~1};
\node[lm,  right=of t1]  (lm2)  {LM~2\\(reason)};
\node[tool, right=of lm2] (t2)   {Tool~2};
\node[lm,  right=of t2]  (lmk)  {LM~$k$\\(respond)};
\node[box, fill=gray!15, right=of lmk] (tts) {TTS};

% Arrows with latency labels
\draw[arr] (asr)  -- (lm1)  node[midway, above, annot]{$\lambda_{asr}$};
\draw[arr] (lm1)  -- (t1)   node[midway, above, annot]{$\lambda_{lm}$};
\draw[arr] (t1)   -- (lm2)  node[midway, above, annot]{$\lambda_{tool}$};
\draw[arr] (lm2)  -- (t2)   node[midway, above, annot]{$\lambda_{lm}$};
\draw[arr] (t2)   -- (lmk)  node[midway, above, annot]{$\cdots$};
\draw[arr] (lmk)  -- (tts)  node[midway, above, annot]{$\lambda_{tts}$};

% Brace below showing total silence gap and irreducible floor
\draw[decorate, decoration={brace, amplitude=8pt, mirror}, thick, color=black!60]
  ([yshift=-0.5cm]asr.south west) -- ([yshift=-0.5cm]tts.south east)
  node[midway, below=0.4cm, annot, align=center]
  {$\Lambda_k = \lambda_{asr} + k\,\lambda_{lm} + (k-1)\,\lambda_{tool} + \lambda_{tts}$\\[2pt]
   \textit{Floor:}\;
   $\Lambda_k^* = \lambda_{asr}^{irr} + k\,\lambda_{lm}^{irr} + \lambda_{tts}^{irr}$};

\end{tikzpicture}
%
  }
  \caption{The $k$-step ReAct agent pipeline.
           Each reasoning step (blue) adds a full LM inference latency $\lambda_{lm}$;
           each tool call (orange) adds $\lambda_{tool}$.
           The total silence gap $\Lambda_k$ grows linearly with $k$,
           making the Cascade Competitiveness Criterion (Definition~\ref{def:ccc})
           increasingly hard to satisfy for tool-intensive agents.}
  \label{fig:react_loop}
\end{figure}

\subsection{Implications for Tool-Using Voice Agents}

Theorem~\ref{thm:react} has several important consequences for the design and deployment
of voice agents built on the ReAct paradigm \citep{yao2023react}.

\paragraph{CCC becomes harder to satisfy.}
For representative values ($\lambda_{lm} = 350$~ms, $\lambda_{tool} = 100$~ms,
$\lambda_{asr} = 350$~ms, $\lambda_{tts} = 200$~ms),
the silence gap for a two-step agent ($k = 2$) is already:
\[
  \Lambda_2 = 350 + 2(350) + 100 + 200 = 1350~\text{ms},
\]
nearly $7\times$ the human turn-taking threshold $\tau_{human} = 200$~ms
established by \citet{levinson2015timing} and \citet{stivers2009universals}.
Each additional reasoning step adds at least $\lambda_{lm} + \lambda_{tool} \approx 450$~ms
to the silence gap, pushing competitiveness further out of reach.

\paragraph{Chain-of-thought reasoning compounds the effect.}
When the agent uses chain-of-thought reasoning \citep{wei2023chainofthought},
each LM step generates a longer token sequence before producing the tool call
or response token.
This increases $\lambda_{lm}$ while improving response quality,
creating a fundamental quality--latency trade-off within each reasoning step.
The trade-off is multiplicative in $k$: an agent running $k$ reasoning steps
with chain-of-thought pays the overhead $k$ times.

\paragraph{Tool call latency is largely irreducible.}
Remote API tool calls involve a network round-trip that cannot be eliminated.
For tools requiring database queries or external service calls,
$\lambda_{tool}^{irr}$ reflects the minimum service response time,
set by physical network latency and server-side processing.
Only local, in-process tool calls have $\lambda_{tool}^{irr} \approx 0$;
real-world agents using web search, calendar, or CRM tools face positive
$\lambda_{tool}^{irr}$ values on the order of 50--200~ms per call
\citep{langgraph2024}.

\paragraph{The irreducible floor grows linearly with \texorpdfstring{$k$}{k}.}
Unlike the single-step case, where $\Lambda^* = 380$~ms in our running example,
a $k$-step agent has $\Lambda_k^* = 380 + (k-1) \cdot \lambda_{lm}^{irr}$ (assuming in-process tools).
For $k = 3$ and $\lambda_{lm}^{irr} = 100$~ms, this gives
$\Lambda_3^* = 380 + 2 \times 100 = 580$~ms,
nearly $3\times$ the human threshold.
No engineering lever can push below this floor without changing the pipeline
architecture: architectural modifications such as speculative execution or
parallel tool dispatch (discussed in Section~\ref{sec:design_levers}) address
the reducible terms but leave $\Lambda_k^*$ intact.
